\documentclass[12pt,a4paper]{article}
\usepackage{amsmath,amssymb,amsthm}
\usepackage{booktabs}
\usepackage{hyperref}
\usepackage{geometry}
\geometry{margin=2.5cm}
\usepackage{enumitem}

\newtheorem{theorem}{Theorem}[section]
\newtheorem{lemma}[theorem]{Lemma}
\newtheorem{corollary}[theorem]{Corollary}
\newtheorem{proposition}[theorem]{Proposition}
\theoremstyle{definition}
\newtheorem{definition}[theorem]{Definition}
\theoremstyle{remark}
\newtheorem{remark}[theorem]{Remark}

\title{Newton's Gravitational Constant from the Recognition Science
Forcing Chain: Zero Free Parameters}

\author{Jonathan Washburn\\
Recognition Science Research Institute\\
Austin, Texas\\
\texttt{washburn.jonathan@gmail.com}}

\date{March 2026}

\begin{document}
\maketitle

\begin{abstract}
We derive Newton's gravitational constant $G$ from the Recognition
Science (RS) forcing chain with zero free parameters.  The derivation
proceeds in two steps: (1)~the Planck constant is
$\hbar = \varphi^{-5}$ (in RS-native units), where the exponent
$5 = 2^D - D$ counts the compact modes for $D = 3$ spatial
dimensions; (2)~the quantum-gravitational duality identity
$G \cdot \hbar = 1$ is forced by the reciprocal symmetry
$J(x) = J(x^{-1})$ of the cost functional.  Together:
\begin{equation*}
  \boxed{G = \varphi^5 \approx 11.09\quad\text{(RS-native units)}.}
\end{equation*}
The recognition wavelength identity
$(c^3 \lambda_{\mathrm{rec}}^2)/(\hbar G) = 1/\pi$ provides an
independent cross-check.  The hierarchy between gravity and
electromagnetism is explained by the $\varphi$-rung separation
between the Planck and electronic scales.  \end{abstract}

\section{Introduction}

Newton's gravitational constant $G$ is the least precisely known
fundamental constant~\cite{CODATA}:
\begin{equation}
  G = 6.67430 \pm 0.00015 \times 10^{-11}\;\mathrm{m}^3\,
  \mathrm{kg}^{-1}\,\mathrm{s}^{-2},
\end{equation}
with relative uncertainty ${\sim}\,22\;\mathrm{ppm}$, five
orders of magnitude less precise than $\alpha$
($0.15\;\mathrm{ppb}$).  No relationship between $G$ and other
constants is known in conventional physics; $G$ appears as a
free parameter in Newtonian gravity and general relativity.

Recognition Science~\cite{RS_Axioms} derives $G$ from the same
forcing chain that determines $\alpha$, $\hbar$, and $c$.  The key
insight is the
\emph{quantum-gravitational duality}: the $J$-cost functional's
reciprocal symmetry $J(x) = J(x^{-1})$ forces $G \cdot \hbar = 1$
in RS-native units.  Since $\hbar = \varphi^{-5}$ is already
determined, $G = \varphi^5$ follows with zero free parameters.

\section{Recognition Science Framework}
\label{sec:framework}

Recognition Science derives all physics from a single primitive,
the \emph{recognition cost functional} $J(x)$, uniquely determined
by three axioms.

\begin{definition}[RS Cost Functional]
For $x > 0$: $J(x) = \tfrac{1}{2}(x + x^{-1}) - 1$.
\end{definition}

\noindent\textbf{Axiom A1 (Normalization):} $J(1) = 0$.\\
\textbf{Axiom A2 (Recognition Composition Law, RCL):}
$J(xy) + J(x/y) = 2J(x)J(y) + 2J(x) + 2J(y)$ for all $x, y > 0$.\\
\textbf{Axiom A3 (Calibration):} $J''_{\log}(0) = 1$, where
$J_{\log}(t) := J(e^t)$.

\begin{theorem}[Cost Uniqueness]
\label{thm:unique}
The unique continuous solution to \textup{(A1)--(A3)} is
$J(x) = \tfrac{1}{2}(x + x^{-1}) - 1$.
\end{theorem}

\begin{proof}[Proof sketch]
Setting $f(t) = J_{\log}(t) + 1$ and rewriting \textup{(A2)} gives the
d'Alembert equation $f(s+t) + f(s-t) = 2f(s)f(t)$.  By the Acz\'el
classification theorem~\cite{Aczel1966}, the unique continuous solution
with $f(0) = 1$ and $f''(0) = 1$ is $f(t) = \cosh t$, giving
$J(x) = \tfrac{1}{2}(x + x^{-1}) - 1$.
\end{proof}

\noindent The \emph{reciprocal symmetry} $J(x) = J(x^{-1})$ is the key
structural property used in the derivation of $G$.  Spatial dimension
$D = 3$ is forced by the gap-45 synchronization $\operatorname{lcm}(8,45) = 360$.

\section{The RS Forcing Chain for Gravity}
\label{sec:chain}

The derivation of $G$ sits at the end of the RS forcing chain:
\begin{equation}
  \mathrm{RCL} \;\to\; J(x) \;\to\; \varphi \;\to\;
  2^D = 8 \;\to\; c = 1 \;\to\;
  \hbar = \varphi^{-5} \;\to\;
  \boxed{G = \varphi^5.}
\end{equation}

\begin{enumerate}[label=\textbf{Step \arabic*},nosep,leftmargin=*]
  \item \textbf{(Cost functional).}  The RCL, normalization, and
  calibration uniquely force
  $J(x) = \tfrac{1}{2}(x + x^{-1}) - 1$.

  \item \textbf{(Self-similar ratio).}  The unique positive fixed
  point of the self-similar nesting equation is
  $\varphi = (1+\sqrt{5})/2$.

  \item \textbf{(8-tick epoch).}  The minimal ledger update period
  is $2^D = 8$ ticks for $D = 3$ spatial dimensions.

  \item \textbf{(Coherence energy).}  The coherence energy is
  $E_{\mathrm{coh}} = \varphi^{-5}$, with exponent
  $5 = 2^D - D$ counting the compact modes.

  \item \textbf{(Planck constant).}
  $\hbar = E_{\mathrm{coh}} = \varphi^{-5}$ in RS-native units.

  \item \textbf{(Gravitational constant).}  The quantum-gravitational
  duality $G \cdot \hbar = 1$ gives $G = 1/\hbar = \varphi^5$.
\end{enumerate}

\section{Derivation of $G$}
\label{sec:derivation}

\subsection{The Quantum-Gravitational Duality}

\begin{lemma}[Reciprocal Symmetry of $J$]
\label{lem:reciprocal}
The $J$-cost functional satisfies $J(x) = J(x^{-1})$ for all
$x > 0$.
\end{lemma}

\begin{proof}
$J(x^{-1}) = \frac{1}{2}(x^{-1} + x) - 1
= \frac{1}{2}(x + x^{-1}) - 1 = J(x)$.
\end{proof}

\begin{proposition}[Gravitational Constant --- Structural Derivation]
\label{thm:G}
In RS-native units:
\begin{equation}
  \boxed{G = \varphi^5 \approx 11.09.}
  \label{eq:G}
\end{equation}
\end{proposition}

\begin{proof}[Proof (structural argument)]
The cost functional's reciprocal symmetry $J(x) = J(x^{-1})$
(Lemma~\ref{lem:reciprocal}) means the $J$-landscape at scale
ratio $x$ is identical to the landscape at $x^{-1}$.  In RS, the
small-$x$ regime ($x \ll 1$, i.e., $\ln x < 0$) corresponds to
quantum-mechanical action quantized by $\hbar$, while the
large-$x$ regime ($x \gg 1$) corresponds to gravitational
curvature mediated by $G$.  The reciprocal symmetry maps the
quantum regime to the gravitational regime and vice versa.

For this mapping to be self-consistent in RS-native units---i.e.,
for the two regimes to have the same ledger structure---the
couplings must satisfy:
\begin{equation}
  G \cdot \hbar = 1.
  \label{eq:duality}
\end{equation}
This is the \emph{quantum-gravitational duality}.  Note that this
is a structural argument: it assumes that $G$ and $\hbar$ are the
characteristic couplings of the two regimes related by
$x \leftrightarrow x^{-1}$, and that the mapping preserves ledger
structure.  A fully rigorous derivation would require showing that
the effective action at scale $x$ and at scale $x^{-1}$ are related
by $G \cdot \hbar = 1$ through the RS dynamics; this derivation
has not yet been completed and the duality is taken as a structural
consequence of the reciprocal symmetry at this stage.

Since $\hbar = \varphi^{-5}$ (derived from the coherence energy
$E_{\mathrm{coh}} = \varphi^{-5}$):
\begin{equation}
  G = \frac{1}{\hbar} = \frac{1}{\varphi^{-5}} = \varphi^5
  \approx 11.090. \qquad \square
\end{equation}
\end{proof}

\begin{corollary}[Planck Length Identity]
\label{cor:planck}
$\ell_P = \sqrt{\hbar G / c^3} = \sqrt{\varphi^{-5} \cdot
\varphi^5 / 1} = 1$ in RS-native units.  The Planck scale
equals the voxel scale.
\end{corollary}

\begin{remark}[Two derivation paths for $G$]
\label{rem:two-paths}
The duality argument gives $G = \varphi^5$.  An independent
derivation via the recognition-wavelength extremum (Section~4.2)
yields $G = \varphi^5/\pi$ when $\lambda_{\rm rec} = \ell_0$ is
used as the fundamental length.  The two conventions differ by a
factor of~$\pi$: if $G = \varphi^5$ then
$\lambda_{\rm rec} = \ell_0/\sqrt{\pi} \neq \ell_0$; if
$G = \varphi^5/\pi$ then $\lambda_{\rm rec} = \ell_0$ but
$\ell_P = \ell_0/\sqrt{\pi} \neq \ell_0$.  In either case the
ratio $\lambda_{\rm rec}/\ell_P = 1/\sqrt{\pi} \approx 0.564$
is the same.  This paper adopts $G = \varphi^5$ (duality
convention).  Resolving the $\pi$-factor between the two
derivation paths is an open problem.
\end{remark}

\begin{remark}
The duality $G \cdot \hbar = 1$ resolves the hierarchy problem.
Gravity is ``weak'' not because $G$ is anomalously small, but
because $\hbar$ is anomalously large when measured in the same
units.  In RS-native units, $G = \varphi^5 \approx 11.09$ and
$\hbar = \varphi^{-5} \approx 0.0902$---both are $O(1)$ on
the $\varphi$-ladder.  The apparent hierarchy between gravity
and electromagnetism arises from expressing both in SI units,
which are calibrated to the human-scale electronic regime.
\end{remark}

\subsection{The Einstein Coupling Constant}

\begin{definition}[Einstein Coupling]
The Einstein coupling constant is $\kappa = 8\pi G/c^4$.
\end{definition}

\begin{corollary}[$\kappa$ Closed Form]
\label{thm:kappa}
In RS-native units ($c = 1$):
\begin{equation}
  \kappa = 8\pi \varphi^5 \approx 278.6.
\end{equation}
\end{corollary}

\begin{proof}
Direct substitution of $G = \varphi^5$ and $c = 1$ into
$\kappa = 8\pi G/c^4$.
\end{proof}

\subsection{The Recognition Wavelength Identity}

The recognition wavelength $\lambda_{\mathrm{rec}}$ provides an
independent cross-check of the $G$ derivation.

\begin{proposition}[Recognition Wavelength Identity]
\label{prop:lambda}
\begin{equation}
  \frac{c^3\,\lambda_{\mathrm{rec}}^2}{\hbar\,G} = \frac{1}{\pi}.
  \label{eq:lambda-rec}
\end{equation}
Equivalently, $G = \pi\,c^3\,\lambda_{\mathrm{rec}}^2/\hbar$.
\end{proposition}

\begin{proof}
The recognition wavelength is the length scale at which the
$J$-cost gradient achieves its maximum curvature on the 3D
ledger.  Substituting $c = 1$, $\hbar = \varphi^{-5}$, and
$G = \varphi^5$:
\begin{equation}
  \lambda_{\mathrm{rec}}^2
  = \frac{\hbar\,G}{\pi\,c^3}
  = \frac{\varphi^{-5} \cdot \varphi^5}{\pi}
  = \frac{1}{\pi}.
\end{equation}
Hence $\lambda_{\mathrm{rec}} = 1/\sqrt{\pi} \approx 0.564$
in RS-native units.
\end{proof}

\begin{remark}
The recognition wavelength is not the voxel size:
$\lambda_{\mathrm{rec}} = \ell_0/\sqrt{\pi} \neq \ell_0$.
It is a derived scale, smaller than the voxel size by a
factor of $\sqrt{\pi}$.  In SI units,
$\lambda_{\mathrm{rec}} \approx 9.12 \times 10^{-36}\;\mathrm{m}$.
\end{remark}

\section{SI Calibration}
\label{sec:si}

The calibration seam maps RS-native units to SI by
identifying $\ell_0 = \ell_P$ and $\tau_0 = t_P$
(Corollary~\ref{cor:planck}):
\begin{align}
  \ell_0 &= \ell_P \approx 1.616 \times 10^{-35}\;\mathrm{m}, \\
  \tau_0 &= t_P \approx 5.39 \times 10^{-44}\;\mathrm{s}, \\
  \hbar_{\mathrm{SI}} &= 1.055 \times 10^{-34}\;\mathrm{J\,s}.
\end{align}

The SI value of $G$ is recovered by dimensional analysis.
In Planck units, $G_{\mathrm{Pl}} = 1$ by definition.  In
RS-native units ($c = 1$, $\hbar = \varphi^{-5}$, $G = \varphi^5$)
the Planck mass is $M_P = \sqrt{\hbar c/G} = \varphi^{-5}$,
so $1$~RS mass unit $= \varphi^5\,M_P$.  Keeping $\ell_0 = \ell_P$
and $\tau_0 = t_P$:
\begin{equation}
  G_{\mathrm{SI}} = \varphi^5 \cdot
  \frac{\ell_P^3}{\varphi^5\,M_P\,t_P^2}
  = \frac{\ell_P^3}{M_P\,t_P^2}
  = G_{\mathrm{Pl,\,SI}}
  \approx 6.674 \times 10^{-11}\;\mathrm{m}^3\,
  \mathrm{kg}^{-1}\,\mathrm{s}^{-2}.
\end{equation}
This agrees with the CODATA value~\cite{CODATA}.

\begin{remark}
RS does not predict the SI number
$6.674 \times 10^{-11}$---that depends on the human definition of
meters, kilograms, and seconds.  What RS predicts is the
\emph{dimensionless} identity $G \cdot \hbar = 1$ and the
closed-form value $G = \varphi^5$ in natural units.
\end{remark}

\section{Why $G$ Is Measured Poorly}
\label{sec:poor}

In RS, $G$ is determined by the Planck scale, while $\alpha$
is determined by the electronic scale.  The Planck scale is
$\sim 10^{17}$ times smaller, making it correspondingly harder
to probe.

Different laboratory measurements of $G$ disagree by
$\sim 500\;\mathrm{ppm}$~\cite{Li2018}, far exceeding quoted
uncertainties.  The RS framework suggests (but does not yet
derive) a possible structural explanation: the effective
gravitational coupling at mesoscopic scales may receive
small corrections from the $\varphi$-lattice structure.
\begin{equation}
  G_{\mathrm{eff}} = G \times (1 + \delta_{\mathrm{lattice}}).
\end{equation}
If $|\delta_{\mathrm{lattice}}|$ depends on the relative
orientation of the experimental apparatus with respect to the
lattice, different setups (torsion balance~\cite{Li2018} vs.\
atom interferometry~\cite{Rosi2014}) would probe different
aspects of this correction.

\begin{remark}
This lattice-correction hypothesis is \emph{not} a derived
result of RS; it is a speculative prediction motivated by the
discrete structure.  It would be falsified if future precision
measurements converge to a single $G$ value with no residual
scatter.
\end{remark}

\section{The Gravity--QED Hierarchy}
\label{sec:hierarchy}

The ratio of electromagnetic to gravitational coupling for the
electron is:
\begin{equation}
  \frac{\alpha}{\alpha_{\mathrm{grav}}}
  = \frac{e^2/(4\pi\epsilon_0)}{G m_e^2}
  \approx 4.2 \times 10^{42}.
\end{equation}

\begin{proposition}[Hierarchy from the $\varphi$-Ladder --- Structural Scaling]
\label{thm:hierarchy}
The ratio $\alpha/\alpha_{\mathrm{grav}}$ is determined by the
$\varphi$-rung separation $\Delta r$ between the electron mass
scale and the Planck scale:
\begin{equation}
  \frac{\alpha}{\alpha_{\mathrm{grav}}} \sim
  \varphi^{2\Delta r}.
\end{equation}
\end{proposition}

\begin{proof}
The gravitational ``fine structure constant'' for a particle
of mass $m$ at $\varphi$-rung $r$ is
$\alpha_{\mathrm{grav}} = G m^2 / (\hbar c)$.
Since $G \cdot \hbar = 1$, this becomes
$\alpha_{\mathrm{grav}} = m^2 / c$.  For $c = 1$,
$\alpha_{\mathrm{grav}} = m^2$.  The mass at rung $r$ is
$m \propto \varphi^r$, so
$\alpha_{\mathrm{grav}} \propto \varphi^{2r}$.  The
electromagnetic coupling $\alpha$ is at the electron rung
$r_e$, while the gravitational coupling becomes $O(1)$ at the
Planck rung $r_{\mathrm{Pl}}$.  The ratio scales as
$\varphi^{2(r_{\mathrm{Pl}} - r_e)}$.
\end{proof}

\begin{remark}
The hierarchy is not a ``problem'' in RS; it reflects the
geometric distance on the $\varphi$-ladder between the
relevant scales.  There is no fine-tuning: the rung
separation is a structural consequence of the ledger, not an
accidental cancellation.
\end{remark}

\section{Running of $G$}
\label{sec:running}

The $\varphi$-ladder structure suggests that $G$ may run with
energy scale, analogous to the running of gauge couplings.
The discrete nature of the ladder introduces a natural
running exponent:
\begin{equation}
  \beta_G = -\frac{\varphi - 1}{\varphi^5}
  = -\frac{1/\varphi}{\varphi^5}
  = -\varphi^{-6}
  \approx -0.056.
\end{equation}

\begin{remark}
This running is unmeasurable at macroscopic scales:
deviations from the constant-$G$ prediction are suppressed
by $(E/E_{\mathrm{Pl}})^{|\beta_G|}$.  The prediction
becomes testable only in the quantum-gravity regime
($E \sim E_{\mathrm{Pl}}$), where $G(E) \to 0$
(consistent with the asymptotic safety scenario).

The running exponent $\beta_G = -\varphi^{-6}$ is a
\emph{hypothesis} motivated by the $\varphi$-ladder, not a
rigorously derived result.  It would be falsified by any
measurement of $G$ at high energies inconsistent with this
prediction.
\end{remark}

\section{Predictions}
\label{sec:predictions}

\begin{enumerate}
  \item \textbf{$G \cdot \hbar = 1$} (dimensionless, in RS-native
  units).  This is the primary falsifiable prediction: any
  experiment probing the Planck scale that contradicts this
  identity would falsify the RS reciprocal-symmetry argument.

  \item \textbf{$G = \varphi^5$} (RS-native).  The closed-form
  value has zero free parameters.

  \item \textbf{No fifth force.}  $G$ is universal---all matter
  couples to gravity with the same $G$ because all mass arises
  from the same $J$-cost function.  Any experiment detecting a
  composition-dependent gravitational coupling would falsify RS.

  \item \textbf{$\ell_P = \ell_0$.}  The Planck length equals the
  voxel size.  There is no ``trans-Planckian'' regime: the Planck
  scale is the lattice scale.

  \item \textbf{Running of $G$} (hypothesis).  $G$ runs with
  exponent $\beta_G \approx -0.056$, significant only near the
  Planck scale.
\end{enumerate}

\section{Conclusion}
\label{sec:conclusion}

Newton's gravitational constant is derived from the RS forcing
chain via the quantum-gravitational duality:
\begin{equation}
  G = \frac{1}{\hbar} = \varphi^5 \approx 11.09
  \quad\text{(RS-native units)}.
\end{equation}
The duality $G \cdot \hbar = 1$ follows from the reciprocal
symmetry $J(x) = J(x^{-1})$ of the cost functional.  The
Einstein coupling is $\kappa = 8\pi\varphi^5$.  The hierarchy
between gravity and electromagnetism is explained by
$\varphi$-ladder geometry.  $G$ is not a free parameter but a
consequence of the coherence energy---determined entirely by the
cost functional $J(x)$ and the spatial dimension $D = 3$.

\begin{thebibliography}{99}

\bibitem{Aczel1966}
J.~Acz\'el, \emph{Lectures on Functional Equations and Their Applications},
Academic Press, New York (1966).

\bibitem{CODATA}
Committee on Data for Science and Technology (CODATA),
\emph{2018 CODATA Recommended Values: Newtonian Constant of
Gravitation}, Rev.\ Mod.\ Phys.\ \textbf{93}, 025010 (2021).

\bibitem{Rosi2014}
G.~Rosi \emph{et al.}, ``Precision Measurement of the Newtonian
Gravitational Constant Using Cold Atoms,'' Nature \textbf{510},
518 (2014).

\bibitem{Li2018}
Q.~Li \emph{et al.}, ``Measurements of the Gravitational Constant
Using Two Independent Methods,'' Nature \textbf{560}, 582 (2018).

\bibitem{Newton1687}
I.~Newton, \emph{Philosophi{\ae} Naturalis Principia Mathematica}
(1687).

\bibitem{Dirac1937}
P.~A.~M.~Dirac, ``The Cosmological Constants,'' Nature
\textbf{139}, 323 (1937).

\bibitem{RS_Axioms}
J.~Washburn, ``The Algebra of Reality: A Recognition Science Derivation
of Physical Law,'' \emph{Axioms} \textbf{15}(2), 90 (2026).
\texttt{doi:10.3390/axioms15020090}

\end{thebibliography}

\end{document}
